% HAND-WRITTEN fragment -- not generated by maestro2tex, but placed by it via the
% `order` list in configs/anatop_tb_tia.json.
% Numbers come from tab_tia_amp, tab_tia_drive, tab_tia_tran and tab_tia_step.

\subsubsection{The amplifier}

Measured in unity feedback into \SI{5}{\pico\farad}, the amplifier is unremarkable and
comfortable: \SIrange{84.1}{85.8}{\decibel} of DC gain, a \SIrange{6.2}{7.6}{\mega\hertz}
unity-gain frequency, and \SI{88}{\degree} of phase margin at every corner. It is
overcompensated for this load, which costs bandwidth and buys an unconditionally clean
step response.

The one asymmetry worth carrying forward is in the output stage. Sourcing collapses at
\SIrange{13.6}{15.6}{\micro\ampere} --- the fixed current in M5 --- while sinking is not
resolved by the \SI{40}{\micro\ampere} sweep, and the same ratio appears large-signal as
a falling slew rate \num{1.7} times the rising one. At the maximum gain setting this
does not bite: full scale through \SI{560}{\kilo\ohm} is \SI{\pm 2.23}{\micro\ampere},
a factor of six inside the sourcing limit. At the \emph{minimum} gain setting it does.
The design target of \SI{\pm 33}{\micro\ampere} at $R_f = \SI{35}{\kilo\ohm}$ requires
the output to pass \SI{33}{\micro\ampere} in both directions, which is more than twice
the sourcing limit measured here. Half of that range --- the half in which the output
must source current through $R_f$ --- is not deliverable. This is an open item: it needs either a larger M5, a rethink of
the class-A output stage, or a reduction of the full-scale current at low gain. It has
not been confirmed in the transimpedance bench itself, because no low-gain setting was
simulated.
